%!TEX root =  tfg.tex
\chapter{Planificación}

\begin{quotation}[Maestro de guerra y escritor]{Sun Tzu (544--496 a.C.)}
Divide y vencerás -El arte de la guerra\end{quotation}

\begin{abstract}
En este capítulo se exponen las divisiones en sprint planificadas para la realización del proyecto.
\end{abstract}

\section{Resumen temporal del proyecto}

\begin{table*}[htb]
	\centering
	\begin{coolTable}{ll}{2}
{Resumen del proyecto}
	\textbf{Fecha de inicio}&13/10/2025\\
	\textbf{Fecha de fin}&16/05/2026\\
	\textbf{Periodicidad de las revisiones}&3 semanas\\
	\textbf{Carga de trabajo semanal}&12 horas\\
	\textbf{Horas totales previstas}&285 horas\\
	\textbf{Horas finales}&293 horas\\
	\end{coolTable}
	\caption{Tabla resumen de tiempos y planificación}
\end{table*}

\section{Planificación inicial}

Aquí un desglose de las iteraciones, comienzo y fin de cada una:

\begin{table*}[htb]
	\centering
	\begin{coolTable}{ll}{2}
{Resumen de iteraciones}
	\textbf{Iteración 1}&13/10/25 a 13/03/26\\ 
	\textbf{Iteración 2}&13/03/26 a 03/04/26\\ 
	\textbf{Iteración 3}&04/04/26 a 24/04/26\\
	\textbf{Iteración 4}&25/04/26 a 08/05/26\\
	\end{coolTable}
	\caption{Planificación temporal de iteraciones}
\end{table*}

La primera iteración cubre un período más largo que las siguientes: el proyecto arranca en octubre de 2025, pero la documentación no comienza hasta febrero de 2026. El trabajo de implementación realizado durante ese período no se descarta, sino que se incorpora en esta primera iteración. Las iteraciones siguientes se planifican con una duración de tres semanas cada una, lo que permite mantener un ritmo constante de trabajo y dedicar revisiones periódicas con el tutor al progreso del proyecto.


